\section{需要的是解，而不是逆矩阵}
\label{sec:08-Numerical-Analysis/07-Numerical-Methods-for-ML/03}

一段算多元正态对数似然的代码，三行，和教科书上的公式逐字对应：

\[
\log p(x)=-\tfrac12(x-\mu)^{\trans}\Sigma^{-1}(x-\mu)-\tfrac12\log\det\Sigma-\tfrac{n}{2}\log(2\pi).
\]

程序照着写：先 \texttt{inv} 求出 \(\Sigma^{-1}\)，再算二次型，再 \texttt{det} 求行列式、取对数。维数 50 时一切正常。换成 500 维，\texttt{det} 返回 \texttt{0.0}，对数是 \(-\infty\)，似然全变成 \texttt{-inf}。

行列式是 \(n\) 个特征值的连乘。特征值都在 \(0.1\) 附近时，\(\det\Sigma\approx10^{-500}\)——这就是上一节那个下溢，换了身衣服而已。而真正要的那个数 \(\log\det\Sigma\approx-1151\)，普通得不能再普通。第二处毛病更隐蔽：那个 \(\Sigma^{-1}\) 从头到尾只被用了一次，用来乘一个向量。为了这一次乘法，程序造了一个有 \(250{,}000\) 个元素的矩阵，其中绝大多数结果最终没被读过。

\subsection{\(A^{-1}b\) 是一个等式，不是一条指令}

推导里写 \(x=A^{-1}b\)，说的是``\(x\) 是那个满足 \(Ax=b\) 的向量''。它陈述了一个关系，没有规定怎么把这个向量算出来。可一旦逐字翻译成代码，就变成了先造出整个逆映射、再把它作用到一个向量上。

这么做要付三笔账。第一笔是算术：对一般稠密方阵，带主元 LU 分解约 \(\tfrac23n^3\) 次运算，之后两次三角回代各 \(n^2\)；而显式求逆本身就要做一次分解，再解 \(n\) 个单位向量右端，总量约 \(2n^3\)，然后还得再花 \(2n^2\) 做一次矩阵向量乘。三倍的代价，换回一个用完就扔的中间结果。

\begin{center}
\begin{tikzpicture}[x=1cm,y=1cm,font=\small,>=stealth]
  \draw[fill=black!6,draw=black!50] (0,2.5) rectangle (1.5,3.4);
  \node[font=\footnotesize] at (0.75,2.95) {\(A,b\)};
  \draw[fill=red!12,draw=black!50] (2.3,2.5) rectangle (6.1,3.4);
  \node[font=\footnotesize,align=center] at (4.2,2.95) {显式求逆 \(A^{-1}\)\\\(\approx2n^3\)};
  \draw[fill=red!12,draw=black!50] (6.9,2.5) rectangle (10.2,3.4);
  \node[font=\footnotesize,align=center] at (8.55,2.95) {与 \(b\) 相乘\\\(2n^2\)};
  \draw[fill=black!6,draw=black!50] (11.0,2.5) rectangle (13.2,3.4);
  \node[font=\footnotesize] at (12.1,2.95) {\(\hat x\)};
  \draw[->,black!60] (1.55,2.95) -- (2.25,2.95);
  \draw[->,black!60] (6.15,2.95) -- (6.85,2.95);
  \draw[->,black!60] (10.25,2.95) -- (10.95,2.95);
  \draw[fill=black!6,draw=black!50] (0,0.8) rectangle (1.5,1.7);
  \node[font=\footnotesize] at (0.75,1.25) {\(A,b\)};
  \draw[fill=green!16,draw=black!50] (2.3,0.8) rectangle (6.1,1.7);
  \node[font=\footnotesize,align=center] at (4.2,1.25) {分解 \(A=LU\)\\\(\approx\tfrac23n^3\)，保带宽};
  \draw[fill=green!16,draw=black!50] (6.9,0.8) rectangle (10.2,1.7);
  \node[font=\footnotesize,align=center] at (8.55,1.25) {两次三角回代\\\(2n^2\)};
  \draw[fill=black!6,draw=black!50] (11.0,0.8) rectangle (13.2,1.7);
  \node[font=\footnotesize] at (12.1,1.25) {\(\hat x\)};
  \draw[->,black!60] (1.55,1.25) -- (2.25,1.25);
  \draw[->,black!60] (6.15,1.25) -- (6.85,1.25);
  \draw[->,black!60] (10.25,1.25) -- (10.95,1.25);
  \node[font=\footnotesize,black!75] at (6.2,2.15) {残差界里多带一个 \(\kappa(A)\)：\(\hat x\) 不是任何邻近系统的精确解};
  \node[font=\footnotesize,black!75] at (6.2,0.45) {后向稳定：残差 \(O(u)\)，前向误差 \(O(\kappa u)\)};
\end{tikzpicture}

\emph{同一个 \(Ax=b\)，两条路径。上面一条多花三倍算术、把稀疏或带状结构抹成稠密矩阵，换回来的答案连残差都不一定小。}
\end{center}

图里下面那行小字是第二笔账，也是最容易被忽略的一笔。带主元的 LU 加回代是后向稳定的：算出的 \(\hat x\) 是某个 \((A+\Delta A)\hat x=b\) 的精确解，\(\norm{\Delta A}/\norm{A}\) 是 \(O(u)\) 量级，于是残差小到机器精度。先求逆再相乘则不是——求逆过程中的舍入进了 \(A^{-1}\) 的元素，再被矩阵乘法搅在一起，最终残差界里会多出一个 \(\kappa(A)\) 因子。病态时两条路径的前向误差都是 \(O(\kappa u)\) 量级，看上去打平，但上面那条连``残差小''这个可以直接测出来的保证都丢了，出问题时没有抓手。这两个指标的分工见\hyperref[sec:08-Numerical-Analysis/01-Floating-Point-and-Error/04]{``前向误差、后向误差与稳定算法''}。

第三笔账是结构。稀疏矩阵的逆几乎总是稠密的：一个 \(10^6\) 阶三对角系统，直接解是 \(O(n)\) 的事，几毫秒；它的逆有 \(10^{12}\) 个元素，先不说算多久，根本存不下。带状、三角、Toeplitz、Kronecker 积——这些结构都会在求逆那一步被抹平。分解把结构留住了：带宽为 \(p\) 的矩阵，其 LU 因子的带宽还是 \(p\)。更大的稀疏问题连分解都不做，共轭梯度和 GMRES 只需要一个``矩阵乘向量''的接口，矩阵本身甚至不必显式存在，见\hyperref[sec:08-Numerical-Analysis/05-Numerical-Linear-Algebra/04]{``大型稀疏系统与 Krylov 迭代''}。

方法的选择跟着``最终要什么''走：一般方阵解方程用带主元 LU；对称正定用 Cholesky；满列秩最小二乘用 QR；需要辨认数值秩用 SVD；同一个矩阵配许多右端，分解一次、回代多次。没有哪一个永远最好，但``不显式求逆''是一个几乎不会错的默认起点。这几种分解的完整取舍在\hyperref[sec:08-Numerical-Analysis/05-Numerical-Linear-Algebra/01]{``线性方程为何分解而不求逆''}。

\subsection{一次 Cholesky，交出三样东西}

回到开头那段代码。协方差矩阵对称正定，Cholesky 分解 \(\Sigma=LL^{\trans}\) 给出下三角的 \(L\)，代价约 \(\tfrac13n^3\)，是 LU 的一半。有了它，公式里的三样东西全都不必碰逆矩阵。

解 \(\alpha=\Sigma^{-1}y\)，就是依次解 \(L z=y\) 与 \(L^{\trans}\alpha=z\)，两次三角回代，各 \(n^2\)。

二次型更省一步。把 \(u=x-\mu\) 代进去，

\[
u^{\trans}\Sigma^{-1}u=u^{\trans}(LL^{\trans})^{-1}u=(L^{-1}u)^{\trans}(L^{-1}u)=\norm{z}_2^2,
\qquad Lz=u.
\]

一次三角回代加一次内积就完了，连 \(\alpha\) 都不用求。

\(\log\det\) 那一项的补救和上一节是同一招。\(\det\Sigma=\det(L)\det(L^{\trans})=(\prod_i L_{ii})^2\)，取对数把连乘换成连加：

\[
\log\det\Sigma=2\sum_i\log L_{ii}.
\]

五百个 \(\log L_{ii}\) 相加不会下溢，不管它们各自多小。会溢出的是那个从没被要求过的中间量 \(\det\Sigma\)——只要不去构造它，问题就不存在。

条件要对一遍。这套做法要求 \(\Sigma\) 严格正定；核方法和高斯过程里的核矩阵常常在数值上处于边缘正定，Cholesky 会在某个对角元处遇到非正数而失败。惯常做法是加一点抖动 \(\Sigma+\varepsilon I\)，但要清楚这一步改的是问题而不是算法——\(\varepsilon\) 是在给模型加观测噪声，它的大小该由建模判断，不该由``让分解跑通''来定。另外，\(L_{ii}\) 中出现极小值本身就是个信号：它说明 \(\kappa(\Sigma)\) 很大，后面所有基于这次分解的结果都得按\hyperref[sec:08-Numerical-Analysis/01-Floating-Point-and-Error/03]{``条件数衡量问题本身的敏感性''}那把尺子打折。

\subsection{正规方程是同一处误译的另一副面孔}

最小二乘 \(\min_x\norm{Ax-b}_2\) 的驻点条件是正规方程 \(A^{\trans}Ax=A^{\trans}b\)，逐字翻译就得到

\[
x=(A^{\trans}A)^{-1}A^{\trans}b.
\]

这里除了求逆，还多了一处更贵的错误：形成 \(A^{\trans}A\) 会把条件数平方。由 \(A=U\Sigma V^{\trans}\) 得 \(A^{\trans}A=V\Sigma^2V^{\trans}\)，奇异值逐个平方，于是 \(\kappa_2(A^{\trans}A)=\kappa_2(A)^2\)。一个 \(\kappa_2(A)=10^8\) 的设计矩阵——特征之间存在中等程度的共线性时很常见——平方之后是 \(10^{16}\)，双精度下一位有效数字都不剩；而直接对 \(A\) 做 QR 分解，误差只按 \(\kappa_2(A)\) 走，答案有八位可信。

改法不是换个函数名，是换一条推导路径：\(A=QR\)，代入后 \(Q\) 的正交性把问题化成 \(Rx=Q^{\trans}b\)，一次三角回代解决，全程没有出现 \(A^{\trans}A\)。矩阵接近秩亏时再往前一步用 SVD，它会明确告诉你哪些方向在数据里几乎无法分辨。这两条路见\hyperref[sec:08-Numerical-Analysis/05-Numerical-Linear-Algebra/02]{``QR 与最小二乘的稳定解法''}与\hyperref[sec:08-Numerical-Analysis/05-Numerical-Linear-Algebra/06]{``SVD、低秩与数值秩''}。

\subsection{求解器治不好病态}

有一点需要说清楚：不求逆是省钱和保结构，不是消除敏感性。若 \(\kappa(A)\) 本来就大，数据里一点扰动就会让精确解大幅移动，这是问题的属性，换任何函数调用都改不掉。

残差在这时最容易被误读。\(r=b-A\hat x\) 很小，移项即得 \(A\hat x=b-r\)，它只说明 \(\hat x\) 精确地解了一个右端被改动过一点的系统；前向误差还要乘上条件数才算得出来。一个近乎奇异的系统，完全可以同时给出很小的残差和很大的解误差。判断时得把残差、缩放、条件数估计和``这个解在业务上是否可能''放在一起看，任何一项单独都不够。

混合精度迭代改进是个有用的折中：低精度做分解，高精度算残差，再解一次修正量。它有效的前提同样写在条件数里，\(\kappa u\) 太大时迭代改进也收不回来——它是一个利用误差结构的算法，不是无条件的加速开关。

\subsection{隐式层的反传，还是解方程}

即便某一层在数学上就写成 \(x=A^{-1}b\)，对它求导也没有理由构造逆。对 \(Ax=b\) 取微分得 \((dA)x+A\,dx=db\)，即

\[
A\,dx=db-(dA)x,
\]

前向敏感度仍是一个以同一个 \(A\) 为系数的线性系统——分解可以直接复用。反向也一样：设上游给出对 \(x\) 的梯度 \(g\)，先解伴随方程 \(A^{\trans}\lambda=g\)，则由

\[
g^{\trans}dx=\lambda^{\trans}A\,dx=\lambda^{\trans}\big(db-(dA)x\big)
\]

可以直接读出对 \(b\) 的梯度是 \(\lambda\)，对 \(A\) 的梯度是 \(-\lambda x^{\trans}\)。这正是上一节所说的``隐式模型的反传只需解一个伴随线性系统''的具体样子：既没有显式 Jacobian，也没有逆矩阵，只多解了一个方程组，而且系数矩阵还是原来那个。可微优化层、Kalman 滤波、可微物理里的反传，走的都是这条路。

\subsection{什么时候逆本身就是要的输出}

``不要求逆''不是禁令。有些任务确实以逆矩阵为输出：广义线性模型要报告参数的协方差矩阵，也就是 Fisher 信息阵的逆的全部元素；实验设计要看 \((X^{\trans}X)^{-1}\) 的谱；留一交叉验证的闭式公式要用到帽子矩阵的对角。这些场合里，完整的逆就是答案，构造它天经地义。

要问的从来不是``能不能求逆''，而是``最终交付的到底是什么''：是一个解，是若干右端，是一个二次型，是一个 \(\log\det\)，还是逆本身。前四种情况占了绝大多数，而它们都有比求逆更快、更稳、更省内存的走法。哪怕真的只要逆矩阵的几个对角元，在大规模问题上，专门的估计方法也仍然远胜于先把整个逆造出来再取几个数。

这一章的三篇讲的是同一件事的三副面孔。概率连乘该换到对数域，导数该交给计算图，逆矩阵该换成解方程——三次改写都没有动数学，动的是通往同一个数的路径。后面的凸优化、统计推断和深度学习里，公式会越来越长，而``这是在陈述一个等式，还是在描述一条计算路径''这个问题，值得每次都问一遍。
